% 卡齐米日·库拉托夫斯基（Kazimierz Kuratowski）（综述）
% license CCBYSA3
% type Wiki

卡齐米日·库拉托夫斯基（Kazimierz Kuratowski，波兰语发音：[kaˈʑimjɛʂ kuraˈtɔfskʲi]；1896年2月2日－1980年6月18日）是一位波兰数学家和逻辑学家。他是华沙数学学派的主要代表人物之一，曾任华沙大学教授以及波兰科学院数学研究所（IM PAN）的研究员。1946年至1953年间，他担任波兰数学学会会长。

他以在集合论、拓扑学、测度论和图论方面的贡献而著称。以他名字命名的一些著名数学概念包括：库拉托夫斯基定理、库拉托夫斯基闭包公理、库拉托夫斯基–佐恩引理以及库拉托夫斯基交集定理。
\subsection{生平与事业}
\subsubsection{早年生活}
卡齐米日·库拉托夫斯基于1896年2月2日出生在华沙（当时属于受俄罗斯帝国控制的波兰会议王国）。他的父亲马雷克·库拉托夫是一名律师，母亲是罗莎·卡热夫斯卡。他在以保卫将军帕韦乌·赫沙诺夫斯基（Paweł Chrzanowski）命名的华沙中学完成中等教育。1913年，他进入苏格兰的格拉斯哥大学学习工程学，部分原因是他不愿在俄语教学体系下学习——当时在波兰禁止使用波兰语授课【需要引用】。他仅学习了一年，第一次世界大战的爆发使进一步的学业中断。1915年，俄军撤出华沙，华沙大学重新开学，并以波兰语作为教学语言。同年，库拉托夫斯基重返大学，这次主修数学。1921年，他在新成立的第二波兰共和国获得博士学位。
\subsubsection{博士论文}
1921年秋，库拉托夫斯基因其开创性的研究工作获得博士学位。他的论文由两部分组成。第一部分致力于通过闭包公理建立拓扑学的公理化体系。这一部分（1922年以略作修改的形式再版）被数百篇科学论文引用过【1】。

库拉托夫斯基论文的第二部分研究的是“两点间不可约连续体”的问题。这一主题最初是齐格蒙特·亚尼谢夫斯基（Zygmunt Janiszewski）在其法语博士论文中的研究对象。由于亚尼谢夫斯基已去世，库拉托夫斯基的论文导师改由斯特凡·马祖凯维奇（Stefan Mazurkiewicz）担任。库拉托夫斯基的研究解决了比利时数学家夏尔-让·艾蒂安·古斯塔夫·尼古拉·德·拉瓦莱-普桑男爵（Charles-Jean Étienne Gustave Nicolas, Baron de la Vallée Poussin）在集合论中提出的若干问题。
\subsubsection{二战前的学术生涯}
两年后的1923年，库拉托夫斯基被任命为华沙大学数学系副教授。1927年，他受聘为利沃夫理工大学（Lwów Polytechnic）的数学正教授，并一直担任该校数学系主任至1933年。库拉托夫斯基还曾两度出任该系的院长。1929年，他成为华沙科学学会的成员。

虽然库拉托夫斯基与利沃夫数学学派的许多学者，如斯特凡·巴拿赫（Stefan Banach）和斯坦尼斯瓦夫·乌拉姆（Stanisław Ulam），以及以苏格兰咖啡馆（Scottish Café）为中心的一群数学家有着密切的学术往来，但他始终与华沙保持紧密联系。1934年，库拉托夫斯基离开利沃夫返回华沙，那是在著名的《苏格兰问题集》（*Scottish Book*）开始编写（1935年）之前，因此他并未向其中贡献任何问题。然而，他确实与巴拿赫在测度论的重要问题上有过密切合作【2】【3】。

1934年，他被任命为华沙大学教授。一年后，库拉托夫斯基被任命为数学系主任。1936年至1939年间，他担任波兰科学与应用科学委员会数学委员会秘书。
\subsubsection{战时与战后}
第二次世界大战期间，由于德国占领当局禁止波兰人接受高等教育，库拉托夫斯基在华沙的地下大学秘密授课。

1945年2月，华沙大学重新开学后，库拉托夫斯基开始在该校讲授课程。同年，他成为波兰学习科学院（Polish Academy of Learning）院士。1946年，他被任命为华沙大学数学系副主任，1949年起，他当选为华沙科学学会副会长。1952年，他成为波兰科学院（Polish Academy of Sciences）院士，并在1957年至1968年间担任该院副院长。

二战后，库拉托夫斯基积极参与波兰科学界的重建。他协助创立了国家数学研究所（State Institute of Mathematics），该机构于1952年并入波兰科学院【4】。1948年至1967年间，库拉托夫斯基担任波兰科学院数学研究所所长，同时长期担任波兰数学学会及国际数学学会的主席。他还在1963年至1966年间担任国际数学联盟（International Mathematical Union）副主席，并于1968年至1980年出任国家数学研究所科学委员会主席。1948年至1980年间，他一直领导该院的拓扑学部门。

他的学生之一是著名逻辑学家兼集合论学者安杰伊·莫斯托夫斯基（Andrzej Mostowski）。
